%% kapa: exam=1 round=21 year=2010 period=1 form=A subject=economics subject_ko=경제원론 seq=22 number=62
%% source: Q-Net 감정평가사 기출문제 (cst00309) artlSeq=5024181 file=extracted/1008/5024181/2010년 제21회 감정평가사 1차시험 기출문제/1교시.hwp
%% conversion: kordoc 4.24.51 -> tools/gen1.py (auto); review=proofread-sonnet; flags=none
\begin{question}{62}
인플레이션에 관한 설명으로 옳은 것은? (단, 다른 조건은 일정하다.)
\begin{choices}
  \item 예상인플레이션율이 증가하면 실질이자율이 상승한다.
  \item 합리적 기대가설에 의하면 예상인플레이션율이 증가할 경우 인플레이션이 심화된다.
  \item 예상치 못한 인플레이션이 발생하면 자연실업률이 하락한다.
  \item 인플레이션율이 예상보다 높으면 고정된 연금으로 생활하는 사람들에게 유리해진다.
  \item 인플레이션을 완전히 예상할 수 있다면 메뉴비용이 발생하지 않는다.
\end{choices}
\end{question}
\answer{2}
